\section{Conclusion}
\label{sec:conclusion}

Three lattice options for a storage-ring search for the electric dipole
moments of light nuclei at JINR have been presented.

The NICA collider can be used as a storage ring for EDM experiments by adding
bypass channels parallel to the original straight sections. Wien filters
installed in the bypasses compensate the MDM spin rotation in the arcs, while
the arcs remain unchanged, so that NICA keeps its collider functions. The most
realistic bypass, with a fully regular straight section, satisfies all optical
requirements, and spin-orbit tracking confirms that the quasi-frozen spin
method can be realized in NICA. The simultaneous control of betatron and spin
chromaticity is under study, and the same lattice allows NICA to be used as an
axion antenna.

The Nuclotron can be modernized for deuteron and proton EDM experiments in
the quasi-frozen spin mode with electrostatic deflectors or Wien filters. If
the Nuclotron serves as the collider booster, its final energy can be lowered
to 2--\SI{3}{GeV}, which frees about \SI{70}{m} of the orbit for
polarization-preserving solenoids and EDM and axion equipment.

Finally, two specialized frozen spin rings have been designed: a purely
electrostatic proton ring with a circumference of \SI{419.1}{m} and a
deuteron ring with combined E+B deflectors with a circumference of
\SI{145.85}{m}. These lattices form the basis for the physical design and cost
estimate of facilities intended for the search for the violation of
fundamental symmetries.
